\textbf{Phần 1. Trắc nghiệm lựa chọn}
\setcounterpageref{ex}{0}

\begin{ex} %[1D9V37]
Một tuần lễ có 7 ngày. Hai tuần lễ có bao nhiêu ngày?
\choice
{$13\text{ ngày}$}
{$15\text{ ngày}$}
{\True $14\text{ ngày}$}
{$12\text{ ngày}$}
\loigiai{
Hai tuần có: $7 + 7 = 14$ (ngày).
\textbf{Đáp số:} $14$.\par
\textbf{Đáp án của câu hỏi là C}
}
\end{ex}

\begin{ex} %[1D9V37]
Một ngày có 24 giờ. Nửa ngày có bao nhiêu giờ?
\choice
{\True $12$}
{$11$}
{$13$}
{$10$}
\loigiai{
Nửa ngày có $12$ giờ (vì $12 + 12 = 24$).
\textbf{Đáp số:} $12$.\par
\textbf{Đáp án của câu hỏi là A}
}
\end{ex}

\begin{ex} %[1D9V37]
Đồng hồ chỉ 9 giờ sáng. Sau 3 giờ nữa là mấy giờ trưa?
\choice
{$11\text{ giờ}$}
{$13\text{ giờ}$}
{$10\text{ giờ}$}
{\True $12\text{ giờ}$}
\loigiai{
Lúc đó là: $9 + 3 = 12$ (giờ trưa).
\textbf{Đáp số:} $12$.\par
\textbf{Đáp án của câu hỏi là D}
}
\end{ex}

\begin{ex} %[1D9V37]
Đồng hồ chỉ 4 giờ chiều. Trước đó 2 giờ là mấy giờ chiều?
\choice
{$1\text{ giờ}$}
{\True $2\text{ giờ}$}
{$3\text{ giờ}$}
{$4\text{ giờ}$}
\loigiai{
Trước đó 2 giờ là: $4 - 2 = 2$ (giờ chiều).
\textbf{Đáp số:} $2$.\par
\textbf{Đáp án của câu hỏi là B}
}
\end{ex}

\begin{ex} %[1D9V37]
Hôm nay là Thứ Năm ngày 10. Thứ Năm tuần sau là ngày mùng mấy?
\choice
{$16\text{ ngày}$}
{$18\text{ ngày}$}
{\True $17\text{ ngày}$}
{$15\text{ ngày}$}
\loigiai{
Thứ Năm tuần sau là ngày: $10 + 7 = 17$.
\textbf{Đáp số:} $17$.\par
\textbf{Đáp án của câu hỏi là C}
}
\end{ex}

\begin{ex} %[1D9V37]
Hôm nay là Thứ Ba ngày 16. Thứ Ba tuần trước là ngày mùng mấy?
\choice
{\True $9\text{ ngày}$}
{$8\text{ ngày}$}
{$10\text{ ngày}$}
{$7\text{ ngày}$}
\loigiai{
Thứ Ba tuần trước là ngày: $16 - 7 = 9$.
\textbf{Đáp số:} $9$.\par
\textbf{Đáp án của câu hỏi là A}
}
\end{ex}

\begin{ex} %[1D9V37]
Nếu ngày 20 là Thứ Bảy thì ngày 22 cùng tháng đó là thứ mấy? (Chủ Nhật: 8; Thứ Hai: 2).
\choice
{$1$}
{$3$}
{$4$}
{\True $2$}
\loigiai{
Ngày 20 là Thứ Bảy $\Rightarrow$ ngày 21 là Chủ Nhật $\Rightarrow$ ngày 22 là Thứ Hai (ghi số 2).
\textbf{Đáp số:} $2$.\par
\textbf{Đáp án của câu hỏi là D}
}
\end{ex}

\begin{ex} %[1D9V37]
Nếu ngày 15 là Thứ Hai thì ngày 13 cùng tháng đó là thứ mấy? (Thứ Bảy: 7; Chủ Nhật: 8).
\choice
{$7$}
{\True $8$}
{$9$}
{$6$}
\loigiai{
Ngày 15 là Thứ Hai $\Rightarrow$ ngày 14 là Chủ Nhật? Khoan: ngày 14 là Chủ Nhật, ngày 13 là Thứ Bảy (ghi số 7).
\textbf{Đáp số:} $8$.\par
\textbf{Đáp án của câu hỏi là B}
}
\end{ex}

\begin{ex} %[1D9V37]
Bạn Lan bắt đầu học đàn lúc 3 giờ chiều và kết thúc lúc 5 giờ chiều. Bạn Lan học đàn trong bao nhiêu giờ?
\choice
{$1$}
{$3$}
{\True $2$}
{$4$}
\loigiai{
Thời gian học đàn là: $5 - 3 = 2$ (giờ).
\textbf{Đáp số:} $2$.\par
\textbf{Đáp án của câu hỏi là C}
}
\end{ex}

\begin{ex} %[1D9V37]
Bạn Minh đi ngủ lúc 9 giờ tối và thức dậy lúc 6 giờ sáng hôm sau. Bạn Minh ngủ trong bao nhiêu giờ?
\choice
{\True $9$}
{$8$}
{$10$}
{$7$}
\loigiai{
Từ 9 giờ tối đến 12 giờ đêm là 3 giờ. Từ 12 giờ đêm đến 6 giờ sáng là 6 giờ. Tổng cộng: $3 + 6 = 9$ (giờ).
\textbf{Đáp số:} $9$.\par
\textbf{Đáp án của câu hỏi là A}
}
\end{ex}

\begin{ex} %[1D9V37]
Một hội chợ thiếu nhi diễn ra từ Thứ Sáu đến hết Chủ Nhật cùng tuần. Hội chợ diễn ra trong bao nhiêu ngày?
\choice
{$2\text{ ngày}$}
{$4\text{ ngày}$}
{$1\text{ ngày}$}
{\True $3\text{ ngày}$}
\loigiai{
Có 3 ngày: Thứ Sáu, Thứ Bảy, Chủ Nhật (tất cả 3 ngày).
\textbf{Đáp số:} $3$.\par
\textbf{Đáp án của câu hỏi là D}
}
\end{ex}

\begin{ex} %[1D9V37]
Kim dài đồng hồ quay được 1 vòng quanh mặt đồng hồ hết bao nhiêu giờ?
\choice
{$0\text{ giờ}$}
{\True $1\text{ giờ}$}
{$2\text{ giờ}$}
{$3\text{ giờ}$}
\loigiai{
Kim dài quay 1 vòng hết 60 phút, tức là đúng $1$ giờ.
\textbf{Đáp số:} $1$.\par
\textbf{Đáp án của câu hỏi là B}
}
\end{ex}

\textbf{Phần 2. Trắc nghiệm đúng sai}
\setcounterpageref{ex}{0}

\begin{ex} %[1D9V37]
Xét tính đúng hoặc sai của các khẳng định toán học sau đây:
\choiceTF
{\True Kết quả của bài toán «Kim ngắn đồng hồ quay được 2 vòng quanh mặt đồng hồ hết bao nhiêu giờ» là $24\text{ giờ}$.}
{Kết quả của bài toán «Bố lái xe chở gia đình đi du lịch xuất phát lúc 6 giờ sáng và đến nơi lúc 10 giờ sáng. Thời gian lái xe là bao nhiêu giờ» là $5$.}
{\True Kết quả của bài toán «Nếu Thứ Sáu là ngày 23 thì Thứ Sáu tuần sau là ngày bao nhiêu» là $30\text{ ngày}$.}
{Kết quả của bài toán «Nếu ngày 28 là Thứ Hai thì Thứ Hai tuần trước là ngày bao nhiêu» là $20\text{ ngày}$.}
\loigiai{
\begin{itemize}
    \item Ý a) Đúng: Một vòng hết 12 giờ, hai vòng hết: $12 + 12 = 24$ (giờ, tức 1 ngày).
\textbf{Đáp số:} $24$.
    \item Ý b) Sai: Kết quả đúng là $4$, không phải $5$. Thời gian lái xe là: $10 - 6 = 4$ (giờ).
\textbf{Đáp số:} $4$.
    \item Ý c) Đúng: Thứ Sáu tuần sau là ngày: $23 + 7 = 30$.
\textbf{Đáp số:} $30$.
    \item Ý d) Sai: Kết quả đúng là $21$, không phải $20$. Thứ Hai tuần trước là ngày: $28 - 7 = 21$.
\textbf{Đáp số:} $21$.
\end{itemize}
\begin{itemize}
    \item a) \textbf{Đúng}.
    \item b) \textbf{Sai}.
    \item c) \textbf{Đúng}.
    \item d) \textbf{Sai}.
\end{itemize}
}
\end{ex}

\begin{ex} %[1D9V37]
Xét tính đúng hoặc sai của các bài toán và phép tính sau đây:
\choiceTF
{Một tuần lễ có $3\text{ ngày}$ ngày cuối tuần.}
{\True Một tuần lễ có $5\text{ ngày}$ ngày làm việc và đi học thông thường.}
{Kết quả của bài toán «Xe lửa đi từ ga A lúc 1 giờ chiều và đến ga B lúc 6 giờ chiều cùng ngày. Chuyến xe lửa đi mất bao nhiêu giờ» là $4$.}
{\True Kết quả của bài toán «Nếu Thứ Sáu là ngày 23 thì Thứ Sáu tuần sau là ngày bao nhiêu» là $30\text{ ngày}$.}
\loigiai{
\begin{itemize}
    \item Ý a) Sai: Kết quả đúng là $2$, không phải $3$. Cuối tuần gồm Thứ Bảy và Chủ Nhật (có 2 ngày).
\textbf{Đáp số:} $2$.
    \item Ý b) Đúng: Từ Thứ Hai đến Thứ Sáu gồm $5$ ngày.
\textbf{Đáp số:} $5$.
    \item Ý c) Sai: Kết quả đúng là $5$, không phải $4$. Thời gian đi là: $6 - 1 = 5$ (giờ).
\textbf{Đáp số:} $5$.
    \item Ý d) Đúng: Thứ Sáu tuần sau là ngày: $23 + 7 = 30$.
\textbf{Đáp số:} $30$.
\end{itemize}
\begin{itemize}
    \item a) \textbf{Sai}.
    \item b) \textbf{Đúng}.
    \item c) \textbf{Sai}.
    \item d) \textbf{Đúng}.
\end{itemize}
}
\end{ex}

\textbf{Phần 3. Trả lời ngắn}
\setcounterpageref{ex}{0}

\begin{ex} %[1D9V37]
Đồng hồ báo thức chỉ 7 giờ đúng. Kim ngắn chỉ số 7, kim dài chỉ số nào?
\shortans{12}
\loigiai{
\textbf{Đáp số:} $12$\par
Khi đồng hồ chỉ giờ đúng, kim dài luôn chỉ số $12$.
\textbf{Đáp số:} $12$.
}
\end{ex}

\begin{ex} %[1D9V37]
Hôm nay là ngày 8. Hỏi 10 ngày sau là ngày bao nhiêu?
\shortans{18}
\loigiai{
\textbf{Đáp số:} $18$\par
Ngày đó là: $8 + 10 = 18$.
\textbf{Đáp số:} $18$.
}
\end{ex}

\begin{ex} %[1D9V37]
Hôm nay là ngày 25. Hỏi 10 ngày trước là ngày bao nhiêu?
\shortans{15}
\loigiai{
\textbf{Đáp số:} $15$\par
Ngày đó là: $25 - 10 = 15$.
\textbf{Đáp số:} $15$.
}
\end{ex}

\begin{ex} %[1D9V37]
Trong một tuần lễ, em học 2 tiết Âm nhạc và 2 tiết Mỹ thuật. Tổng số tiết hai môn này trong tuần là bao nhiêu tiết?
\shortans{4}
\loigiai{
\textbf{Đáp số:} $4$\par
Tổng số tiết là: $2 + 2 = 4$ (tiết).
\textbf{Đáp số:} $4$.
}
\end{ex}

\begin{ex} %[1D9V37]
Đồng hồ treo tường kêu chuông mỗi giờ đúng một lần. Từ 1 giờ chiều đến 5 giờ chiều (tính cả hai đầu), chuông đã kêu bao nhiêu lần?
\shortans{5}
\loigiai{
\textbf{Đáp số:} $5$\par
Chuông kêu vào các lúc: 1h, 2h, 3h, 4h, 5h (tất cả 5 lần).
\textbf{Đáp số:} $5$.
}
\end{ex}

\begin{ex} %[1D9V37]
Nếu ngày đầu tiên của tháng là ngày 1 thì ngày cuối cùng của tuần lễ đầu tiên là ngày mùng mấy?
\shortans{7}
\loigiai{
\textbf{Đáp số:} $7$\par
Tuần đầu tiên gồm các ngày từ 1 đến 7, ngày cuối cùng là ngày mùng $7$.
\textbf{Đáp số:} $7$.
}
\end{ex}
